\documentclass{article}
\usepackage[zh]{homework}

\config{分析\,-\,0}{2026 春季}{12}

\begin{document}


\begin{problem}[7-1]
试证每个一致收敛的有界函数序列一致有界.
\end{problem}

\begin{proof}
  设 \( f_n \rightrightarrows f \), 且 \( f_n \) 对于每一个正整数 \( n \) 有界. 取 \( \varepsilon=1 \), 知存在正整数 \( N \), 使得对每一个 \( n>N \), 均有 \( \sup | f_n-f | <1 \). 而 \( f_n \) 有界, 故 \( f \) 也有界.

  设 \( \sup |f| = M \), 并且设 
  \[ M'=1+\max\left\{ \sup |f_n-f| :1\le n\le N \right\}. \]
  则对于每一个 \( n \), 都有
  \[ \sup |f_n| \le \sup |f| + \sup |f_n-f| <M+M'. \]
  而 \( M+M' \) 是一个常数, 故 \( f_n \) 的上确界是有界的. 因此 \( f_n \) 是一致有界的.
\end{proof}



\begin{problem}[7-2]
如果 $\{f_n\}$ 与 $\{g_n\}$ 在集 $E$ 上一致收敛, 试证: $\{f_n + g_n\}$ 在 $E$ 上一致收敛. 此外, 假如 $\{f_n\}$ 与 $\{g_n\}$ 都是有界函数的序列, 试证 $\{f_n g_n\}$ 在 $E$ 上一致收敛.
\end{problem}

\begin{proof}
  设 \( f_n \rightrightarrows f \), 并且 \( g_n \rightrightarrows g \). 则对于任意的正实数 \( \varepsilon \), 存在 \( N_1, N_2 \), 使得对于任意的 \( n>N_1 \), 有 \( \sup |f_n-f| <\varepsilon/2 \); 对于任意的 \( n>N_2 \), 有 \( \sup |g_n-g| <\varepsilon/2 \). 则对于任意的 \( n>\max\{N_1,N_2\} \), 均有
  \[ \sup|(f_n+g_n)-(f+g)| \le \sup |f_n-f| + \sup |g_n-g| \le \frac{\varepsilon}{2} + \frac{\varepsilon}{2} = \varepsilon. \]

  若 \( f, g \) 均有界, 设 \( \sup\{|f|,|g|\}=M \). 对于任意的 \( 0<\varepsilon <1 \), 存在 \( N_1,N_2 \), 使得对于任意的 \( n>N_1 \), 有 \( \sup |f_n-f| <\varepsilon/(2M+1) \); 对于任意的 \( n>N_2 \), 有 \( \sup |g_n-g| <\varepsilon/(2M+1) \). 则对于任意的 \( n>\max\{N_1,N_2\} \), 均有
  \begin{align*}
    \sup|f_n g_n - fg| 
    &\le \sup |f_n||g_n-g| + \sup |g||f_n-f| \\
    &\le \left(M+\frac{\varepsilon}{2M+1} \right)\cdot \frac{\varepsilon}{2M+1} + M\cdot \frac{\varepsilon}{2M+1} <\varepsilon.
  \end{align*}
  故 \( f_ng_n \rightrightarrows fg \). 
\end{proof}



\begin{problem}[7-3]
作两个序列 $\{f_n\}$, $\{g_n\}$. 要它们在某个集 $E$ 上一致收敛, 但是 $\{f_n g_n\}$ 在 $E$ 上不一致收敛 (当然, $\{f_n g_n\}$ 应在 $E$ 上收敛).
\end{problem}

\begin{solution}
  取 \( E=(0,1) \), \( f_n(x)=g_n(x)=1/x+1/n \). 则 \( f_n \rightrightarrows 1/x \), \( g_n \rightrightarrows 1/x \).
  
  但是 \( f_n g_n \rightrightarrows 1/x^2 \) 却不是一致收敛的, 因为
  \[ f_ng_n-fg = \left( \frac{1}{x}+\frac{1}{n} \right)^2 - \frac{1}{x^2} = \frac{2}{xn} + \frac{1}{n^2} \to+\infty\ (x\to 0^+). \]
\end{solution}



\begin{problem}[7-4]
研究级数
\[
f(x) = \sum_{n=1}^\infty \frac{1}{1 + n^2 x}.
\]
这级数对于 $x$ 的什么值绝对收敛? 它在什么闭区间上一致收敛? 在什么闭区间上它失去一致收敛性? 是否在级数收敛的地方, $f$ 都是连续的? $f$ 是有界的吗?
\end{problem}

\begin{solution}
  对于 \( x\neq 0 \), 当 \( n \) 足够大时, \( (1+n^2x)^{-1} \) 的符号与 \( x \) 的符号相同, 并且 \( (1+n^2x)^{-1} \sim 1/(xn^2) \), 因此 \( f(x) \) 当 \( x\neq 0 \) 时绝对收敛.

  对于闭区间 \( [a,b] \) 满足 \( 0<a<b \), \( (1+n^2x)^{-1} \le 1/(an^2) \), 由 Weierstrass-M 判别法可知 \( f \) 在 \( [a,b] \) 上一致收敛. 同理, \( f \) 在闭区间 \( [a,b] \) 满足 \( a<b<0 \) 上一致收敛. 若 \( 0\in [a,b] \), 那么由 \( f \) 在 \( 0 \) 处发散, 知 \( f \) 在 \( [a,b] \) 上不一致收敛.

  对于任意的 \( x_0 \neq 0 \), 可以找到一个包含 \( x_0 \) 但不包含 \( 0 \) 的闭区间 \( [a,b] \), 从而 \( f \) 在 \( [a,b] \) 上一致收敛, 因此 \( f \) 在 \( x_0 \) 处连续. 

  对于任意的正整数 \( N \), 由
  \[ f\left( \frac{1}{N^2} \right) > \sum_{n=1}^{N} \frac{1}{1+n^2/N^2} > \sum_{n=1}^{N} \frac{1}{2} = \frac{N}{2} \]
  知 \( f \) 是无界的.
\end{solution}



\begin{problem}[7-8]
假设
\[
I(x) =
\begin{cases}
0 & (x \le 0),\\
1 & (x > 0),
\end{cases}
\]
若 $\{x_n\}$ 是 $(a, b)$ 内相异点的序列. 又若 $\sum_{n=1}^\infty \abs{c_n}$ 收敛, 试证级数
\[
f(x) = \sum_{n=1}^\infty c_n I(x - x_n) \quad (a \le x \le b)
\]
一致收敛, 对每个 $x \ne x_n$, $f$ 连续.
\end{problem}

\begin{proof}
  先证明 \( f \) 一致收敛. 由于对于任意的正整数 \( n \), \( \sup |c_nI(x-x_n)| =|c_n| \), 而 \( \sum_{n=1}^{+\infty} c_n \) 一致收敛, 故由 Weierstrass-M 判别法可知 \( f \) 一致收敛.

  再证明 \( f \) 在 \( x\ne x_n \) 处连续. 对于任意的 对于任意的 \( \varepsilon>0 \), 存在 \( N \), 使得对于任意的 \( n>N \), 有 \( |c_n|<\varepsilon \). 并且, 由于 \( x\ne x_n \), 存在 \( \delta>0 \), 使得对于任意的 \( n\le N \), 都有 \( |x-x_n|>\delta \). 则对于任意的 \( y \) 满足 \( |y-x|<\delta \), 都有 \( I(x-x_n)=I(y-x_n) \) (\( n\le N \)), 从而
  \[ |f(y)-f(x)| = \left| \sum_{n=1}^N c_n (I(y-x_n)-I(x-x_n)) + \sum_{n=N+1}^{+\infty} c_n (I(y-x_n)-I(x-x_n)) \right| < \sum_{n=1}^{N}|c_n| < \varepsilon. \]
  故 \( f \) 在 \( x\ne x_n \) 处连续.
\end{proof}



\begin{problem}[7-11]
假设 $\{f_n\}$, $\{g_n\}$ 都是在 $E$ 上定义的, 并且\\
 (a) $\sum f_n$ 的部分和一致有界,\\
 (b) 在 $E$ 上 $g_n \to 0$ 是一致的,\\
 (c) 对于每个 $x \in E$, $g_1(x) \ge g_2(x) \ge g_3(x) \ge \cdots$.\\
试证 $\sum f_n g_n$ 在 $E$ 上一致收敛. 提示: 对照定理 3.42.
\end{problem}

\begin{proof}
  记 \( \sum f_n \) 的部分和为 \( F_n=\sum_{i=1}^{n}f_i \). 则 \( \{F_n\} \) 一致有界, 设 \( \sup |F_n| =M \). 由 Abel 求和, 知
  \[ \sum_{n=1}^{N} f_ng_n 
  = F_ng_n + \sum_{n=1}^{N-1}F_n(g_n-g_{n+1}). \]

  由前面的习题, 知 \( F_ng_n \rightrightarrows 0 \). 并且, 对于每一个 \( x\in E \), 由于 \( g_n(x)-g_{n+1}(x)\ge 0 \), 因此
  \[ \abs{\sum_{n=N_1}^{N_2-1}F_n(x)(g_n(x)-g_{n+1}(x))} 
  \le \sum_{n=N_1}^{N_2-1}\abs{F_n(x)}(g_n(x)-g_{n+1}(x))
  \le M(g_{N_1}(x)-g_{N_2}(x)). \]

  对于任意的 \( \varepsilon>0 \), 由 \( g_n \rightrightarrows 0 \) 与 \( F_ng_n \rightrightarrows 0 \), 知存在 \( N \), 使得对于任意的 \( N<N_1<N_2 \), 均有
  \[ \sup\abs{F_{N_2}g_{N_2}-F_{N_1}g_{N_1}} < \varepsilon, \quad 
  \sup\abs{g_{N_2}-g_{N_1}} < \varepsilon. \]
  则对于任意的 \( N<N_1<N_2 \) 以及任意的 \( x\in E \), 都有
  \begin{align*}
    \abs{\sum_{n=N_1}^{N_2-1} f_n(x)g_n(x) }
    & \le \abs{ F_{N_2}(x)g_{N_2}(x) - F_{N_1}(x)g_{N_1}(x) } 
    + \abs{\sum_{n=N_1}^{N_2-1} F_{n}(x)(g_n(x)-g_{n+1}(x))}  \\
    & \le \varepsilon + M\abs{g_{N_2}(x)-g_{N_1}(x)} < (M+1)\varepsilon.
  \end{align*}

  由于 \( M \) 是与 \( \varepsilon \) 无关的常数, 并且 \( \varepsilon \) 是任意的, 故 \( \sum f_ng_n \) 在 \( E \) 上一致收敛.
\end{proof}



\begin{problem}[7-14]
设 $f$ 是 $\R^1$ 上具有下列性质的连续实函数: $0 \le f(t) \le 1$, 对于每个 $t$, $f(t + 2) = f(t)$, 并且
\[
f(t) =
\begin{cases}
0 & (0 \le t \le 1 / 3),\\
1 & (2 / 3 \le t \le 1).
\end{cases}
\]
令 $\Phi(t) = (x(t), y(t))$, 其中
\[
x(t) = \sum_{n=1}^\infty 2^{-n} f(3^{2n - 1} t), \quad y(t) = \sum_{n=1}^\infty 2^{-n} f(3^{2n} t).
\]
试证 $\Phi$ 连续, 而且 $\Phi$ 把 $I = [0, 1]$ 映满单位正方形 $I^2 \subset \R^2$. 如果已证出这件事, 再证 $\Phi$ 能把 Cantor 集映满单位正方形 $I^2 \subset \R^2$.
\end{problem}

\begin{proof}
  先来证明 \( \Phi \) 是连续函数. 只需证明 \( x(t) \) 与 \( y(t) \) 都是连续函数. 由于对于任意的 \( t \), 均有 \( f(t)\in [0,1] \), 因此 \( \sup\abs{2^{-n}f(3^{2n-1}t)} = 2^{-n} \), 故由 Weierstrass-M 判别法可知 \( x(t) \) 一致收敛. 同理, \( y(t) \) 也一致收敛. 因此 \( x(t) \) 与 \( y(t) \) 都是连续函数, 从而 \( \Phi \) 是连续函数.

  再来证明 \( \Phi \) 可以把 Cantor 集映满单位正方形 $I^2 \subset \R^2$. 对于每一个 \( (x_0,y_0) \in I^2 \), 将它们展开为二进制小数:
  \[ x_0=\sum_{n=1}^{+\infty}2^{-n}a_n, \quad 
  y_0 = \sum_{n=1}^{+\infty}2^{-n}b_n, \quad a_n,b_n\in\{0,1\}. \]
  取
  \[ t_0=\sum_{n=1}^{+\infty}3^{-2n}(2a_n) 
  + \sum_{n=1}^{+\infty}3^{-2n+1}(2b_n). \]

  则
  \begin{align*}
    f(3^{2m-1}t_0) &= f\left( \sum_{n=1}^{+\infty}3^{2m-2n-1}(2a_n) 
  + \sum_{n=1}^{+\infty}3^{2m-2n}(2b_n) \right)\\
  &= f\left( \sum_{n=0}^{+\infty}3^{-2n-1}(2a_{n+m}) 
  + \sum_{n=1}^{+\infty}3^{-2n}(2b_{n+m}) \right)\\
  &=a_m.
  \end{align*}

  同理, \( f(3^{2m}t_0)=b_m. \), 因此有 \( \Phi(t_0)=(x_0,y_0) \). 再注意到 \( t_0 \) 是 Cantor 集中的点 (因为 \( t_0 \) 的三进制小数的每一位均为 \( 0 \) 或 \( 2 \)), 故 \( \Phi \) 可以把 Cantor 集映满单位正方形 $I^2 \subset \R^2$.
\end{proof}



\begin{problem}[7-15]
假设 $f$ 是 $\R^1$ 上的实连续函数, $f_n(t) = f(nt)$, $n=1, 2, 3, \cdots$, 而且 $\{f_n(t)\}$ 在 $[0, 1]$ 上等度连续. 关于 $f$ 你能得出什么结论?
\end{problem}

\begin{solution}
  我们来证明: \( f \) 是常值函数. 否则, 设存在 \( x_1<x_2 \), 并且 \( f(x_1)\neq f(x_2) \). 取正数 \( \varepsilon < \abs{f(x_2)-f(x_1)} \), 则对于任意的正数 \( \delta \), 存在足够大的正整数 \( N \), 使得 \( x_1/N<x_2/N<1 \), 并且 \( (x_2-x_1)/N<\delta \). 因此 \( \abs{f_N(x_1/N)-f_N(x_2/N)} = \abs{f(x_1)-f(x_2)} >\varepsilon \), 并且 \( \abs{x_1/N-x_2/N}<\delta \). 这表明 \( f_n \) 不等度连续, 矛盾. 因此 \( f \) 是常值函数.
\end{solution}



\begin{problem}[7-18]
设 $\{f_n\}$ 是一致有界的函数序列, 这些函数都在 $[a, b]$ 上 Riemann 可积. 令
\[
F_n(x) = \int_a^x f_n(t) \d t \quad (a \le x \le b).
\]
求证存在着子序列 $\{F_{n_k}\}$, 它在 $[a, b]$ 上一致收敛.
\end{problem}

\begin{proof}
  我们先来证明: \( \{F_n(x)\}  \) 在 \( [a,b] \) 上等度连续. 设 \( M=\sup\{f_n(x):x\in[a,b], n\in\N_+\} \). 则对于任意的 \( \varepsilon>0 \), 以及任意的 \( x,y\in[a,b] \) 满足 \( |x-y|<\varepsilon/M \), 都有
  \[ |F_n(x)-F_n(y)| = \left| \int_y^x f_n(t) \d t \right| \le M|x-y| < \varepsilon. \]

  我们再来证明: \( \{F_n(x)\} \) 在 \( [a,b] \) 上逐点有界. 注意到对于所有的正整数 \( n \), 均有 \( F_n(a)=0 \). 根据上面的估计, 对于任意的 \( x\in[a,b] \), 均有
  \[ \abs{F_n(x)}\le M(x-a). \]
  从而 \( \{F_n(x)\} \) 在 \( [a,b] \) 上逐点有界.

  由 Arzelà-Ascoli 定理, 存在子序列 \( \{F_{n_k}\} \) 在 \( [a,b] \) 上一致收敛.
\end{proof}



\begin{problem}[7-20]
如果 $f$ 在 $[0, 1]$ 上连续, 而且
\[
\int_0^1 f(x) x^n \d x = 0 \quad (n = 0, 1, 2, \cdots),
\]
试证在 $[0, 1]$ 上 $f(x) = 0$. 提示: $f$ 与任何多项式之积的积分是零. 用 Weierstrass 定理证明
\[
\int_0^1 f^2(x) \d x = 0.
\]
\end{problem}

\begin{proof}
  由于 \( f \) 在 \( [0,1] \) 上连续, 故 \( f \), 在 \( [0,1] \) 上有界. 设 \( M=\sup\{|f(x)|:x\in[0,1]\} \). 由 Weierstrass 定理, 知对于任意的正实数 \( \varepsilon \), 存在多项式 \( P \), 使得对于任意的 \( x\in[0,1] \), 都有 \( |f(x)-P(x)|<\varepsilon/M \). 则
  \[ 0= \int_0^1 f(x) P(x) \d x = \int_0^1 f^2(x) \d x + \int_0^1 (f(x)-P(x))f(x) \d x. \]

  而
  \[ \abs{\int_{0}^{1} (f(x)-P(x))f(x) \d x} \le M \int_0^1 |f(x)-P(x)| \d x < \varepsilon. \]
  从而
  \[ \abs{\int_0^1 f^2(x) \d x} < \varepsilon. \]
  根据 \( \varepsilon \) 的任意性, \( f(x) = 0 \) 在 \( [0,1] \) 上成立.
\end{proof}



\begin{problem}[7-22]
假定 $f \in \mathcal{R}[a,b]$, 求证有多项式 $P_n$ 合于
\[
\lim_{n\to\infty} \int_a^b \abs{f - P_n}^2 \d x = 0.
\]
(与第 6 章习题 12 对照.)
\end{problem}

\begin{proof}
由于 $f$ 在 $[a,b]$ 上 Riemann 可积, 存在连续函数 $g$ 使得
\[
\int_a^b |f(x)-g(x)|^2 \d x < \frac{\varepsilon}{2}.
\]

对连续函数 $g$, 由 Weierstrass 逼近定理, 存在多项式 $P$ 使得对一切 $x\in[a,b]$ 有 $|g(x)-P(x)|<\sqrt{\varepsilon/(2(b-a))}$, 从而
\[
\int_a^b |g(x)-P(x)|^2 \d x < \frac{\varepsilon}{2}.
\]
于是
\[
\int_a^b |f(x)-P(x)|^2 \d x
\le 2\int_a^b |f(x)-g(x)|^2 \d x + 2\int_a^b |g(x)-P(x)|^2 \d x
< 2\cdot\frac{\varepsilon}{2} + 2\cdot\frac{\varepsilon}{2} = 2\varepsilon.
\]
因此对任意 $\varepsilon>0$ 存在多项式 $P$ 使得 $\int_a^b |f-P|^2 < 2\varepsilon$。  
取 $\varepsilon_n = 1/(2n)$, 则存在多项式 $P_n$ 满足 $\int_a^b |f-P_n|^2 < 1/n$, 从而
\[
\lim_{n\to\infty} \int_a^b |f-P_n|^2 = 0.
\]
\end{proof}



\begin{problem}[7-23]
设 $P_0 = 0$, 对于 $n = 0, 1, 2, \cdots$ 定义
\[
P_{n+1}(x) = P_n(x) + \frac{x^2 - P_n^2(x)}{2}.
\]
试证在 $[-1, 1]$ 上 $\lim_{n \to \infty} P_n(x) = \abs{x}$ 是一致的.
\end{problem}

\begin{proof}
  注意到 \( P_0 \) 是偶函数. 而若 \( P_n \) 是偶函数, 则 \( P_{n+1} \) 也是偶函数. 因此 \( P_n \) 对所有的 \( n \) 都是偶函数. 因此只需证明 \( \{P_n(x)\} \) 在 \( [0,1] \) 上一致收敛于 \( x \). 

  注意到 \( P_n(0)=0 \) 对于所有的非负整数 \( n \) 成立, 因此多项式 \( x-P_n(x) \) 有一个根为 \( 0 \), 从而可以写成 \( x-P_n(x) = xQ_n(x) \). 则
  \[ Q_{n+1}(x) = (1-x)Q_n(x)+\frac{1}{2}xQ_n(x)^2 = Q_n(x)\left( 1-x+\frac{1}{2}xQ_n(x) \right). \]

  先来证明: 对于所有的非负整数 \( n \), 均有 \( Q_n(x)\le 1 \). 命题对 \( n=0 \) 成立; 假设命题对 \( n \) 成立, 则 \( Q_{n+1}(x)=Q_{n+1}(1-x+xQ_n(x)/2)\le 1-x+x/2\le 1 \). 因此 \( Q_n(x)\le 1 \) 对所有的非负整数 \( n \) 成立.

  因此, 对于所有的 \( n\ge 0 \), \( Q_{n+1}(x)\le Q_n(x)(1-x+x/2)=(1-x/2)Q_n(x) \). 结合 \( Q_0(x)=1 \), 知 \( Q_n(x)\le (1-x/2)^n \). 从而
  \[ x - P_n(x) = xQ_n(x) \le x \left( 1-\frac{x}{2} \right)^n. \]

  结合 \( x\le P_n(x) \) (因为 \( P_{n+1}(x)=-(1-P_n(x)^2)/2+(1+x^2)/2 \le -(1-x)^2/2+(1+x^2)/2 = x \)), 知
  \[ \sup\{\abs{x-P_n(x)}: x\in[0,1]\} = \sup\left\{ x\left( 1-\frac{x}{2} \right)^n: x\in[0,1] \right\}\le \frac{2n^n}{(n+1)^{n+1}}\to 0.\ (n\to\infty) \]
  因此 \( P_n(x) \rightrightarrows x \) (\( x\in[0,1] \)) 成立. 从而 \( P_n(x) \rightrightarrows |x| \) (\( x\in[-1,1] \)) 成立.
\end{proof}



\begin{problem}[7-24]
设 $X$ 是以 $d$ 为度量的度量空间. 固定一点 $a \in X$. 给每个 $p \in X$ 指派一个函数 $f_p$,
\[
f_p(x) = d(x, p) - d(x, a) \quad (x \in X).
\]
试证 $\abs{f_p(x)} \le d(a, p)$ 对于一切 $x \in X$ 成立, 从而 $f_p \in \mathscr{C}(X)$. 对于一切 $p, q \in X$, 证明
\[
\lVert f_p - f_q \rVert = d(p, q)
\]
如果 $\Phi(p) = f_p$, 那么 $\Phi$ 是从 $X$ 到 $\Phi(X) \subset \mathscr{C}(X)$ 的等距(保持距离的)映射. 设 $Y$ 是 $\Phi(X)$ 在 $\mathscr{C}(X)$ 里的闭包, 试证 $Y$ 是完备的.
\end{problem}

\begin{proof}
  由三角不等式, 知 \( \abs{f_p(x)} = \abs{d(x,p)-d(x,a)} \le d(p,a) \).

  对于任意的 \( x\in X \), 
  \[ \abs{f_p(x)-f_q(x)} = \abs{d(x,p)-d(x,q)} \le d(p,q). \]
  并且, 注意到 \( \abs{f_p(p)-f_q(p)} = d(p,q) \), 因此 \( \left\|f_p-f_q\right\|=d(p,q) \).

  下面来证明 \( Y=\overline{\Phi(X)} \) 是完备的. 设 \( \tilde{X} \) 是 \( X \) 的完备化. 将 \( \Phi \) 的定义域拓展为 \( \tilde{X} \), 即重新定义 \( \Phi: \tilde{X} \to \mathscr{C}(X) \), \( p\mapsto f_p:=d_{\tilde{X}}(x,p)-d_X(x,a) \).

  我们来证明: \( \Phi(\tilde{X}) \subset \overline{\Phi(X)} \). 对于任意的 \( f_p\in \Phi(\tilde{X}) \), 若 \( p\in X \), 则 \( f_p\in\Phi(X) \); 否则, \( p\in \tilde{X}-X \), 则 \( p \) 的任意邻域 \( N_\delta(p) \) 均含有 \( X \) 中的某一点 \( p' \). 由 \( \left\|f_p-f_{p'}\right\|=d(p,p') \) 知 \( f_{p'}\in N_{\delta}(f_p) \). 由于 \( \delta \) 是任意的, 故 \( f_p\in\overline{\Phi(X)} \).

  对于 \( \overline{\Phi(X)} \) 中的 Cauchy 列 \( \{f_{p_n}\} \), 由于 \( \Phi \) 是等距映射, 故 \( \{p_n\} \) 在 \( X \) 中也是 Cauchy 列, 故收敛于 \( \tilde{X} \) 中的某个点 \( p \), 因此 \( f_{p_n} \to f_p \). 又由于 \( f_p \in \overline{\Phi(X)} \), 故 \( Y=\overline{\Phi(X)} \) 是完备的.
\end{proof}



\begin{problem}
  设 \( \{f_n\} \) 是紧集 \( K \) 上等度连续的函数序列. 
  \begin{enumerate}[label=(\alph*)]
    \item 若 \( \{f_n\} \) 在 \( K \) 上逐点有界, 证明 \( \{f_n\} \) 在 \( K \) 上一致有界.
    \item 若 \( \{f_n\} \) 在 \( K \) 上逐点收敛, 证明 \( \{f_n\} \) 在 \( K \) 上一致收敛.
  \end{enumerate}
\end{problem}

\begin{proof}
  (a) 对于 \( x\in K \), 记 \( g(x) = \sup\{f_n(x):n\ge 1\} \). 由于 \( \{f_n\} \) 在 \( K \) 上逐点有界, 故 \( g \) 对于每一个 \( x\in K \) 均有定义. 我们只需要证明 \( g \) 在 \( K \) 上是有界的. 下面来证明 \( g \) 是一致连续的.

  否则, 存在 \( \varepsilon_0>0 \), 使得对任意的 \( \delta>0 \), 存在 \( x,y\in K \) 使得 \( d(x,y)<\delta \) 但 \( |g(x)-g(y)|\ge\varepsilon_0 \). 由于 \( \{f_n\} \) 是等度连续的, 故存在 \( \delta_0 \), 使得任意的 \( x,y\in K \) 满足 \( d(x,y) < \delta_0 \), 以及任意的正整数 \( n \), 都有 \( |f_n(x)-f_n(y)|<\varepsilon_0/3 \). 
  
  取 \( \delta=\delta_0 \), 知存在 \( x_0,y_0\in K \), 满足 \( d(x_0,y_0)<\delta_0 \) 并且 \( \abs{g(x_0)-g(y_0)} \ge \varepsilon_0 \). 不妨设 \( g(x_0) > g(y_0) \). 由 \( g \) 的定义, 知存在 \( n \), 使得 \( \abs{g(x_0)-f_n(x_0)} < \varepsilon_0/3 \) . 从而
  \[ \abs{g(x_0)-f_n(y_0)} \le \abs{g(x_0)-g(y_0)} + \abs{g(y_0)-f_n(y_0)} <\frac{2\varepsilon_0}{3}. \]
  因此
  \[ f_n(y_0) > g(x_0) - \frac{2\varepsilon_0}{3} \ge g(y_0) + \frac{\varepsilon_0}{3}. \]
  这与 \( f_n(y_0) \le g(y_0) \) 矛盾! 因此 \( g \) 是一致连续的. 由于 \( K \) 是紧集, \( g \) 在 \( K \) 上是有界的. 对于下界可以做类似的讨论. 从而 \( \{f_n\} \) 在 \( K \) 上一致有界.

  (b) 设 \( f_n\to f \). 我们先来证明: \( f \) 是连续的. 取一个固定的点 \( x\in K \). 对于任意的 \( \varepsilon>0 \), 存在 \( \delta>0 \), 使得对于任意的 \( n\in \N_+ \) 以及 \( y\in N_\delta(x) \), 都有 \( \abs{f_n(y)-f_n(x)}<\varepsilon/2 \). 因此, 对于任意的 \( y\in N_\delta(x) \), 以及足够大的 \( n\in \N_+ \), 有
  \[ \abs{f_n(y)-f_n(x)} < \varepsilon/2,\quad \abs{f_n(x)-f(x)} < \varepsilon/2 \implies \abs{f_n(y)-f(y)} < \varepsilon. \]
  因为此式对于任意充分大的 \( n \) 均成立, 因此 \( \abs{f(y)-f(x)}<\varepsilon \). 从而 \( f \) 在 \( x \) 处连续. 由于 \( x \) 是 \( K \) 中的任意点, 故 \( f \) 在 \( K \) 上连续.
  
  若 \( \{f_n\} \) 不是一致收敛的, 则存在 \( \varepsilon_0 \), 使得存在一个 \( \N_+ \) 的无穷子集 \( A \), 满足对于任意的 \( n\in A \), 存在 \( x_n\in K \), 使得 \( \abs{f_n(x_n)-f(x_n)} >\varepsilon_0 \). 由于 \( K \) 是紧集, 故 \( \{x_n\} \) 存在极限点, 设为 \( x_0 \). 由于 \( f \) 连续, 存在 \( \delta_1 \), 使得对于任意的 \( y\in N_{\delta_1}(x_0) \), 都有 \( \abs{f(y)-f(x_0)}<\varepsilon_0/3 \).

  由于 \( \{f_n\} \) 是等度连续的, 存在 \( \delta_2 \), 使得对于任意的 \( x,y\in K \) 满足 \( d(x,y)<\delta_2 \), 以及任意的正整数 \( n \), 都有 \( |f_n(x)-f_n(y)|<\varepsilon_0/3 \). 取 \( \delta_0=\min\{\delta_1,\delta_2\} \). 对于充分大的 \( n\in A \), 有 \( d(x,x_0)<\delta_0 \), 因此 
  \[ \abs{f_n(x_0)-f(x_0)} \ge \abs{f_n(x_n)-f(x_n)} - \abs{f_n(x_n)-f_n(x_0)} - \abs{f_n(x_0)-f(x_0)} > \varepsilon_0 - \frac{\varepsilon_0}{3} - \frac{\varepsilon_0}{3} = \frac{\varepsilon_0}{3}. \]

  因此, 对于充分大的 \( n\in A \), \( \abs{f_n(x_0)-f(x_0)} > \varepsilon_0/3 \), 这与 \( \lim_{n\to\infty} f_n(x_0) = f(x_0) \) 矛盾! 因此 \( \{f_n\} \) 在 \( K \) 上一致收敛.
\end{proof}



\begin{problem}
  设 \( f\in C([0,1]) \). 定义 Bernstein 多项式
  \[ B_nf(x) = \sum_{k=0}^{n}f\left( \frac{k}{n} \right)\binom{n}{k}x^k(1-x)^{n-k}, x\in [0,1]. \]
  证明多项式序列 \( \{B_nf(x)\} \) 一致收敛于 \( f \). 
\end{problem}

\begin{proof}
  由于 \( f \) 在 \( [0,1] \) 上连续, 故 \( f \) 是有界且一致连续的. 因此对于任意的 \( \varepsilon>0 \), 存在 \( \delta>0 \), 使得对于任意的 \( x,y\in[0,1] \) 满足 \( |x-y|<\delta \), 都有 \( |f(x)-f(y)|<\varepsilon \). 设 \( M=\sup\{\abs{f(x)}:x\in[0,1]\} \).

  固定 \( x\in[0,1] \). 对于充分大的 \( n \), 有
  \begin{align*}
    &\abs{B_nf(x)-f(x)}
    = \abs{ -f(x) + \sum_{k=0}^{n} f\left( \frac{k}{n} \right) \binom{n}{k}x^k(1-x)^{n-k} } \\
    \le& \abs{ -f(x) + \sum_{\substack{ 0\le k\le n \\ \abs{x-(n/k)}<\delta }}^{}f\left( \frac{k}{n} \right) \binom{n}{k}x^k(1-x)^{n-k} }
    +\abs{ \sum_{\substack{ 0\le k\le n \\ \abs{x-(n/k)}\ge\delta }}^{}f\left( \frac{k}{n} \right) \binom{n}{k}x^k(1-x)^{n-k} }\\
    \le& \abs{f(x)}\left( 1-\sum_{\abs{x-(n/k)}<\delta }^{}\binom{n}{k}x^k(1-x)^{n-k} \right)
    + \sum_{\abs{x-(n/k)}<\delta }^{}\abs{f\left(\frac{k}{n}\right)-f(x) } \binom{n}{k}x^k(1-x)^{n-k}\\
    &+ M \sum_{\abs{x-(n/k)}\ge\delta }^{}\binom{n}{k}x^k(1-x)^{n-k}\\
    \le& \varepsilon + 2M \sum_{\abs{x-(n/k)}\ge\delta }^{} \binom{n}{k}x^k(1-x)^{n-k}.
  \end{align*}

  注意到
  \begin{align*}
    \sum_{\abs{x-(n/k)}\ge\delta }^{} \binom{n}{k}x^k(1-x)^{n-k}
    <& \sum_{\abs{x-(n/k)}\ge\delta }^{} \frac{(x-n/k)^2}{\delta^2}\binom{n}{k}x^k(1-x)^{n-k}
    =\frac{x(1-x)}{n\delta^2}.
  \end{align*}
  因此
  \[ \abs{B_nf(x)-f(x)} \le \varepsilon+ \frac{2M x(1-x)}{n\delta^2}. \]
  由于 \( \varepsilon \) 可以任意小, 并且 \( n \) 可以任意大, 故 \( B_nf(x) \rightrightarrows f(x) \) 在 \( [0,1] \) 上成立.
\end{proof}


\end{document}